\section{Introduction}
\label{sec:intro}

Cosmological anisotropy is often reported through a scalar amplitude, a preferred direction, or a ratio to a reference scale. Such quantities can be useful. Their interpretation nevertheless depends on the tensor information retained by the observation, the statistical law assigned to that information, and the response relating it to a physical field. These are separate questions. A narrow coefficient interval does not resolve an unspecified response, and an accurately propagated uncertainty does not recover information lost before propagation.

Distance catalogues provide a concrete setting. Their angular residuals can contain a calibration dipole, a velocity-dependent vector response, and a quadrupole. A quadrupole has five signed symmetric trace-free (STF) components. Its amplitude alone does not retain its orientation or eigenvalue shape. Moreover, a finite-redshift magnitude quadrupole is an observational coefficient; its conversion to local shear requires a congruence, normalization, and control of the optical approximation. General luminosity-distance cosmography supplies a physical framework for that distinction, but smoothness of a local expansion does not supply a numerical remainder bound at the catalogue redshifts \cite{Heinesen2021,Maartens2024}. Existing angular expansion analyses likewise provide context for apparent anisotropy estimators, rather than a reason to regard a new fit as a new observable in itself \cite{Kalbouneh2026}.

We develop a focused application with two source-bound carriers. The principal carrier is a fifteen-coordinate mean parameter for 274 released distance moduli within a fixed frame of 628 targets. Its normalized row likelihood permits an exact finite-sample confidence construction under explicit Gaussian and fixed-design assumptions. The second carrier is a 36-coordinate apparent Hubble-rate summary from four overlapping windows of 35\,970 Cosmicflows-4 groups \cite{Tully2023}. It retains a full joint covariance and a previously constructed coupled finite measure. We use it to expose how anchors, positive floors, angular weights, and cross-window dependence change functional summaries. The two carriers have different dimensions, units, and source selections. They are not combined into one likelihood.

The statistical ingredients are established. Partial identification separates uncertainty about an identified set from the breadth of that set \cite{ImbensManski2004}. Bounded changes of sampling weights provide sensitivity classes rather than measured selection probabilities \cite{AronowLee2013}. Auxiliary calibration constrains weights through observed totals, with additional assumptions needed to transport outcomes \cite{Deville1992}. Joint-region projection and Gaussian regression pivots are standard constructions. We use these methods to ask what the released observations identify, which functional domains they support, and which distinctions are erased by compression.

Our contribution therefore lies in the combination of an actual selected-data likelihood, preserved tensor and nuisance dependence, source-linked model comparison, and explicit nonidentification witnesses. Three results make this more than a routine propagation exercise. First, a full-rank statistical tensor can coexist with a physical-response fibre that leaves the entire declared admissible shear domain unresolved. Whether this domain is bounded is a separate assumption; unboundedness is asserted only for the formal unrestricted affine extension. Second, the missing portion of a known target frame gives an exact finite-completion bound that remains wide even when observed responses are known. Third, a positive floor can erase all variation in a coupled depth contrast without changing any observed coefficient. Each result identifies an assumption that a scalar plot alone would conceal.

The analysis is retrospective. We reuse fixed rows, coefficient fits, folds, and coupled numerical values; no new observational fit, null simulation, or draw is needed for the results reported here. Where a reference distribution is exact, exactness refers to the specified probability model. It is not a statement that the catalogue selection or covariance law has been empirically established. We retain this qualification when comparing the models and when interpreting the confidence regions.
